\scanpage{392}\mspage{306}

Il faut voir que si $B$ est un ens. de cardinal $\leq\pi_0$, alors le foncteur diagonal

\[
\delta:J\xrightarrow{\deleted{\unsure{\delta}}}J^B \tag*{$(*)$}
\]

est cofinal. Or comme $A$ est grd devant $\pi_0$, \inserted{$\mathrm{card}\,B\leq\pi_0$,} les produits \unsure{indexés} par $B$ commutent aux $\underset{\longrightarrow}{\lim}$ de type $A$ (SGA 4 I 9.8 appliqué à $\mathrm{Cat}$ et à $\pi_0$). Donc la flèche $(*)$ \deleted{\unsure{est la}} s'interprète comme limite inductive des flèches similaires

\[
\delta_r:J_r\longrightarrow J_r^B, \tag*{$(**)$}
\]

lesquelles sont cofinales, donc $\delta$ est cofinal comme limite inductive de foncteurs cofinaux.

\uline{Lemme} 4.13.13.2. Sous les hypothèses de 4.13.13, soient $M$ une catégorie\inserted{ satisfaisant \unsure{à} $\unsure{L^{\rightarrow}_{\pi_0}|L_{\pi_0}}$}, et $(x_j)_{j\in J}$ \inserted{(i.e. $J\longrightarrow M$)} un système inductif \deleted{$J\longrightarrow M$} dans $M$, indexé par $J$. Pour tout $J'\in I$, posons

\[
x_{J'}=\underset{i\in J'}{\underset{\longrightarrow}{\lim}}\,x_j
\]

% archive sheet number: 391
